\documentclass[12pt]{article}
\usepackage[utf8]{inputenc}
\usepackage[T1]{fontenc}
\usepackage{hyperref}
\title{Minz Zhuang Language Resources for K-12 Curriculum Development}
\author{Compiled by Claude Sonnet 5.5 (claude-sonnet-5-5) for Andr\'as Kornai}
\date{}
\begin{document}
\maketitle

\section{Introduction}
Minz Zhuang (ISO 639-3 zgm; Glottocode minz1236) is a Tai language of southern China in the Kra-Dai family. Sources disagree on where it is spoken. Wikipedia places it in the Langheng area of Funing County, Yunnan, and possibly in southwestern Guangxi, while Global Recordings Network places it in Guangxi. Speaker estimates range from about 170,000 to 183,000. Speakers are reported to be bilingual in Yei (Guibian) Zhuang. The searches found no school materials in Minz Zhuang. Very little exists, and this dossier lists only what was found.

\section{Primary Educational Resources}
No primers, textbooks or school programmes in Minz Zhuang were found in the search results. Nothing is known to be taught in schools.

\section{Linguistic and Reference Materials}
\begin{itemize}
\item Wikipedia, ``Min Zhuang language'': overview of classification, location, speaker numbers and alternate names. \url{https://en.wikipedia.org/wiki/Min_Zhuang_language}
\item Wikitongues language page for Minz Zhuang. \url{https://wikitongues.org/languages/zgm/}
\end{itemize}

\section{Religious and Community Resources}
\begin{itemize}
\item Joshua Project, Minz Zhuang in China. Its summary reports that Bible portions were published in 2025, with no complete Bible or New Testament yet. \url{https://joshuaproject.net/people_groups/20323}
\item Joshua Project language page for zgm. \url{https://joshuaproject.net/index.php/languages/zgm}
\end{itemize}

\section{Digital and Audio Resources}
\begin{itemize}
\item Global Recordings Network, ``Zhuang, Minz language'' (audio recordings listing). \url{https://www.globalrecordings.net/en/language/zgm}
\end{itemize}

\section{Curriculum Development Context}
No existing Minz Zhuang educational materials were found. A curriculum would have to start from documentation work: an orthography decision, a lexicon and recorded speech. Because speakers are reportedly bilingual in Yei (Guibian) Zhuang, and Standard Zhuang is the written standard for the macrolanguage, materials in those varieties may serve as a bridge. Local community consultation would be essential.

\end{document}